% !TeX root = ../nc-quadric-classification.tex
% 2026-11-22 / v1.43（証明をBへ移管）
\section{noncubic 境界を横断する有限型効果性}
\label{sec:noncubic-finite-type}

局所環上の構成を一つの有限型近傍へ広げる。
有限 foundation の Lagrangian 表示に全幾何点を認識する有限開条件を課し、
その同じ基底の全 fiber の geometric helix を用いて全 Hom を先に平坦化する。
二つの cubic 開部分と非被約 noncubic 閉部分から Noetherian 性を貼り合わせ、
局所中央被覆を有限型開へ延長する。
一般の Koszul 性の開性や、無限個の開集合の交叉の開性は用いない。
以下の相対結果とその組合せの新規性は未確認である。

\subsection{Lagrangian 開部分の全幾何点の認識}

以下では標数零の代数閉体上で考える。
前節の Grassmann 表示で、第二 summand の $D_L$ による基底変更を行い、
$H=U\oplus U^*$、$U=k^2\otimes k^2\simeq\operatorname{Mat}_2(k)$ と書く。
交代形式は二 summand 間の標準評価から作るものとする。
この基底変更は二次元因子上の同型なので分解テンソルの集合を保つ。
$\Sigma_U\subset\mathbb P(U)$ と $\Sigma_{U^*}\subset\mathbb P(U^*)$ を
二つの Segre 二次曲面とする。

\begin{lemma}[両階数三の Lagrangian 平面の直接正規形]
\label{lem:lagrangian-double-rank-three-recognition}
$L\subset U\oplus U^*$ を Lagrangian 四平面とする。
二つの射影の階数がともに三で、$L\cap U$ と $L\cap U^*$ の各生成元が
二次正方行列として可逆であると仮定する。
このとき $\operatorname{GL}_2\times\operatorname{GL}_2$ の $U$ 上の tensor 作用と
$U^*$ 上の反傾作用、および $L$ の基底変更により、$L$ は
\begin{equation}\label{eq:lagrangian-rank-three-normal-form}
\begin{aligned}
 L=\bigl\langle&(-E_{21},\alpha E_{12}+\gamma E_{21}),\quad
 (E_{12},\beta E_{12}-\alpha E_{21}),\\
 &(-D_0,0),\quad (0,I_2)\bigr\rangle,
 \qquad \alpha^2+\beta\gamma=1.
\end{aligned}
\end{equation}
となる。ここで $D_0=\operatorname{diag}(1,-1)$ である。
対応する有限 foundation は、中央二重平方環
\[
 k[u,v]\langle c,d\rangle/
 \bigl(c^2-uv,\ d^2-(\alpha u+\beta v)(\gamma u-\alpha v)\bigr)
\]
の既知の normal foundation と同型である。
この認識は、候補 foundation 自身の domain 性や $\mathcal H_2$ 所属を仮定しない。
\end{lemma}
\begin{proof}
全証明は相対モジュライ論文 \cite[noncubic有限型効果性]{QuadricModuli2026} に置く。
\end{proof}

\begin{proposition}[非被約閉基底上の両階数三正規形]
\label{prop:lagrangian-relative-double-rank-three}
$R$ を標数零の体上の可換環、$L\subset U_R\oplus U_R^*$ を
Lagrangian な階数四の部分束とする。
二射影の三次小行列式がそれぞれ単位イデアルを生成し、その四次行列式が
$R$ の元としてともに零であると仮定する。
さらに二つの交点線束の局所生成元は、二次正方行列として可逆とする。
このとき $\operatorname{Spec}R$ の \texorpdfstring{\'etale}{etale} 被覆上で、
前補題の式~\eqref{eq:lagrangian-rank-three-normal-form} が成立する。
従ってその有限 foundation は同じ基底上の明示中央二重平方族の
有限 foundation と同型である。基底の被約性は不要である。
\end{proposition}
\begin{proof}
全証明は相対モジュライ論文 \cite[noncubic有限型効果性]{QuadricModuli2026} に置く。
\end{proof}

\paragraph{Cross reshuffling.}
$M\in\operatorname{Hom}(U,U^*)$ を四因子テンソルと見て、その係数を
$M_{ck,ab}$ と書く。cross reshuffling を
\[
 \operatorname{Cr}(M)_{ka,bc}=M_{ck,ab}
\]
で定める。これは中間の隣接二因子を入力とする縮約を表す四次正方行列である。
$M=y\otimes\nu$ が階数一で、$y$ と $\nu$ がそれぞれ二次正方行列として可逆なら、
$\operatorname{Cr}(M)$ は行・列の並替を除いて二つの可逆二次行列の Kronecker 積であり、可逆である。

\begin{proposition}[noncubic 点を含む全幾何点認識の開部分]
\label{prop:lagrangian-geometric-recognition-open}
$L_0$ を前補題の条件を満たす点とし、有限 foundation の Grassmann chart 内にあるとする。
$\operatorname{LG}(4,H)$ における $L_0$ の開近傍 $W$ を取れて、その全幾何点の
有限 foundation は $\mathcal H_2$ 対象の normal foundation となる。
具体的に、$L$ の局所基底を列行列
\[
 \begin{pmatrix}M\\N\end{pmatrix},\qquad M,N\in\operatorname{Mat}_4
\]
で表し、次の有限開条件を課せばよい。
\begin{enumerate}
\item $\mathbb P(L\cap U)\cap\Sigma_U=\varnothing$、
      $\mathbb P(L\cap U^*)\cap\Sigma_{U^*}=\varnothing$。
\item $\operatorname{Cr}(N\operatorname{adj}M)$ と
      $\operatorname{Cr}(M\operatorname{adj}N)$ が可逆である。
\end{enumerate}
二つの射影がともに階数三なら前補題の noncubic 型であり、
少なくとも一方の階数が四なら一方の ruling foundation が geometric cubic 型である。
この命題は幾何点での有限 foundation の認識であり、相対 heart の存在を結論に先取りしない。
\end{proposition}
\begin{proof}
全証明は相対モジュライ論文 \cite[noncubic有限型効果性]{QuadricModuli2026} に置く。
\end{proof}

\begin{corollary}[cubic 二開部分と noncubic 閉部分]
\label{cor:lagrangian-open-closed-recognition}
前命題の $W$ において、二つの行列式を零とする閉部分スキームを
$D=V(\det M,\det N)$ と置く。
$D$ は非被約でもよい。$D$ 上では命題~\ref{prop:lagrangian-relative-double-rank-three} の正規形が
\texorpdfstring{\'etale}{etale} 局所に成立し、対応する有限 foundation は
相対中央二重平方族のものに一致する。
開補集合は $D(\det M)$ と $D(\det N)$ で覆われ、
前者は元の ruling の geometric cubic 表示、後者は交換した ruling の
geometric cubic 表示を持つ。
\end{corollary}
\begin{proof}
全証明は相対モジュライ論文 \cite[noncubic有限型効果性]{QuadricModuli2026} に置く。
\end{proof}

\subsection{全幾何 fiber の有限 foundation 認識からの一様化}

\begin{theorem}[normal foundation の全 fiber 認識による全軌道の一様化]
\label{thm:uniform-orbit-from-normal-fibers}
$R$ を本稿の基礎体 $k$ 上の可換 Noetherian 代数、$\Lambda$ を有限 projective $R$-代数とし、
$\operatorname{Perf}\Lambda$ は相対 smooth proper とする。
$F_0,F_1,F_2,F_3$ を相対 full strong exceptional foundation とし、
$\operatorname{End}(\bigoplus_{a=0}^3F_a)\simeq\Lambda$ の同定を固定する。
実際の可逆 perfect Morita 双加群による自己同値 $\mathsf s$ と同型
\[
 \mathsf s^2\simeq S^{-1}[2],\qquad
 \mathsf sF_0\simeq F_2,\qquad \mathsf sF_1\simeq F_3
\]
を固定する。ここで $S$ は canonical な相対 Serre 関手である。
各幾何点 $x$ において、$\operatorname{Perf}\Lambda_x$ とある体上の
$\HH$ 対象 $(\mathcal C_x,\mathcal O_x)$ の導来圏との同値があり、
$(F_0,F_1,F_2,F_3)_x$ を定理~\ref{thm:normal-helix} の
normal helix の最初の四項に送ると仮定する。
この仮定は各 fiber の有限 foundation に関するものであり、
$R$ 上の heart や相対正規形を先に仮定しない。

\[
 E_{2a}=\mathsf s^aF_0,\qquad E_{2a+1}=\mathsf s^aF_1,
 \qquad C_{ij}=\operatorname{Hom}_{\Lambda}(E_{-j},E_{-i})\quad(i\le j)
\]
と置き、$i>j$ では $C_{ij}=0$ とする。この同じ $R$ 上で次が成立する。
\begin{enumerate}
\item 全 $i\le j$ に対して $\mathbf R\operatorname{Hom}(E_i,E_j)$ は
  次数零の有限 projective $R$-加群に集中し、その階数は
  \[
  h_{2n}=(n+1)^2,\qquad h_{2n+1}=(n+1)(n+2)
  \]
  である。全 Hom と合成は任意の導来基底変更と整合する。
\item 全軌道環
  \[
  B=\bigoplus_{n\ge0}\operatorname{Hom}(F_0,\mathsf s^nF_0)
  \]
  は一次生成・二次・相対 Koszul であり、全 $B_n$ は階数 $(n+1)^2$ の
  有限 projective 加群である。二次関係は階数七の直和因子となり、
  双対 $B^!$ の全次数、左右および双加群 Koszul 解消は任意基底変更と両立する。
\item $C$ の各頂点単純加群には、距離四以内の五頂点 bar 解消があり、
  これは任意基底変更後にも完全である。
  任意の可換 $R$-代数 $R'$ に対し、canonical な比較
  \[
  D(\operatorname{QGr}B_{R'})\simeq D(\operatorname{Mod}\Lambda_{R'}),
  \qquad D(\operatorname{QGr}B_{R'})^c\simeq\operatorname{Perf}\Lambda_{R'}
  \]
  があり、$\pi B$ を $F_0$ に、$(1)$ を最初に選んだ実際の $\mathsf s$ に送る。
  この比較とその Serre 平方同型は、非平坦なものを含む導来基底変更と整合する。
\end{enumerate}
さらに $B$ が左右 Noetherian なら、任意の有限型可換 $R$-代数 $R'$ 上で、
$\operatorname{qgr}B_{R'}$ は Noetherian heart、module-object 基底変更と fppf descent、
一般 coherent Hom の有限性と平坦基底変更、全 $\chi_q$、AZ 豊富性を持つ。
構造対象は平坦かつ exceptional であり、全幾何 fiber は仮定した $\HH$ の heart と
構造対象を保って exact 同値である。
同じ根の相対 AS4 被覆の存在は、この定理の結論とは別に確認する。
\end{theorem}

\begin{proof}
全証明は相対モジュライ論文 \cite[noncubic有限型効果性]{QuadricModuli2026} に置く。
\end{proof}

\begin{remark}[有限型認識で必要な入力の順序]
この定理では、有限 foundation の全幾何 fiber を既知の normal または cubic 型へ
認識する有限型開条件が、全次数の一様性を供給している。
各次数の平坦 locus を無限に交叉する必要はない。
ここでの大域的な結論は有限 projective 性であり、無限個の Hom 束を
一つの開上で同時に自由化するという、不要な強い主張は含めていない。
また $B$ の Koszul 性を使って全 forward Hom を証明してはいない。
全 fiber の geometric helix から Hom を先に平坦化し、その後に相対 Koszul 原理を使う。
Noetherian 性と同じ根の AS4 被覆は、この順序の後に別途供給すべき入力である。
\end{remark}

\subsection{閉部分と二つの cubic 開部分からの Noetherian 性}

\begin{lemma}[標準 noncubic 族の相対 normal tilting]
\label{lem:nc-relative-normal-tilting}
$R$ を可換 Noetherian $k$-代数とし、
$f_1=uv$、$f_2=(\alpha u+\beta v)(\gamma u-\alpha v)$、
$\alpha^2+\beta\gamma=1$ とする。
\[
 B_0=R[u,v]\langle c,d\rangle/(c^2-f_1,d^2-f_2),\qquad
 Z=R[u,v,h],\quad h=cd+dc
\]
と置く。式~\ref{eq:nchh-plane} の行列分解の cokernel を $M_c,M_d$ とすれば、
両者は $Z$-加群として
\[
 M_c\simeq Z^{\oplus2}\oplus Z(-1)^{\oplus2},\qquad
 M_d\simeq Z^{\oplus2}\oplus Z(-1)^{\oplus2}
\]
であり、これらの表示は任意の基底変更と両立する。
特に四項
\[
 T=(\pi B_0,\pi M_c,\pi B_0(1),\pi M_d(1))
\]
は相対 full strong exceptional foundation で、その End 代数は
式~\ref{eq:nchh-plane} の有限 foundation 代数に一致する。
対応する相対 Morita 同値は四対象、合成、相対 Serre 関手を保ち、
任意の Noetherian 基底変更と整合する。
\end{lemma}
\begin{proof}
全証明は相対モジュライ論文 \cite[noncubic有限型効果性]{QuadricModuli2026} に置く。
\end{proof}

\paragraph{四周期からの Hom $\mathbb Z$-代数の比較。}
標識付き四項 foundation の同型から得た相対 Morita 同値を $F$、
両側の反 Serre 関手を $t,t_0$ と書く。
上の比較は自然同型 $\eta:Ft\simeq t_0F$ を与える。
最初の四項の同型 $\varphi_a:F(E_a)\simeq E_a^0$ ($0\le a<4$) と、
選んだ実際の平方同型による四周期を用い、
$i=a+4n$ に対して $\eta$ の反復から
$\varphi_i:F(E_i)\simeq E_i^0$ を定める。負の $n$ には逆同型を使う。
\[
 f\longmapsto\varphi_jF(f)\varphi_i^{-1}
 \qquad(f:E_i\longrightarrow E_j)
\]
は合成を保つ全 Hom $\mathbb Z$-代数の同型であり、任意基底変更と両立する。
この比較に、二つの根そのものを同一視する仮定は要らない。
異なる $\varphi_i$ の選択は頂点ごとの共役を与えるだけなので、
偶数制限の row と column の Noetherian 性、および線型解消の存在は変わらない。



\begin{lemma}[二つの cubic 開部分と noncubic 閉部分からの普遍 Noetherian 性]
\label{lem:lg-two-rulings-universal-noetherian}
$R$ を可換 Noetherian 環とし、有限 LG 認識開上の foundation と
実際の Morita 根 $s$ を取る。四項を $E_0,E_1,E_2,E_3$ と書き、
指定された同型
\[
 sE_0\simeq E_2,\qquad sE_1\simeq E_3,
 \qquad s^2\simeq S^{-1}[2]
\]
を用いて $E_{2n}=s^nE_0$, $E_{2n+1}=s^nE_1$ と置く。
全ての前向き導来 Hom は有限 projective 加群に次数零で集中し、
任意の底変更と両立することを既に確認したとする。
このとき同じ根による偶数軌道環
\[
 C_n=\operatorname{Hom}(E_0,s^nE_0),\qquad
 a*b=s^{\deg b}(a)\circ b
\]
は、任意の可換 Noetherian $R$-代数 $T$ に対して
$C_T=C\otimes_RT$ が左右 Noetherian となる。

ここで有限 LG 認識開について用いる入力は次の三つである。
二つの射影行列を $M,N$ として
\[
 a=\det M,\quad b=\det N,\quad I=(a,b),
 \qquad Z=\operatorname{Spec}(R/I)
\]
と置く。
\begin{enumerate}
\item $D(a)$ では選んだ ruling の helix が geometric cubic 型である。
\item $D(b)$ では反対の ruling の helix が geometric cubic 型である。
\item $Z$ 上では、非被約構造を含めて、foundation が \'etale 局所的に
  判別式の可逆な二形式による標準 noncubic foundation と同型である。
\end{enumerate}
これらは追加の全次数条件ではなく、有限 LG 認識で確認する条件である。
\end{lemma}

\begin{proof}
全証明は相対モジュライ論文 \cite[noncubic有限型効果性]{QuadricModuli2026} に置く。
\end{proof}

\subsection{同じ根に対する中央 AS4 被覆の延長}

\begin{lemma}[局所中央被覆の有限型近傍への延長]
\label{lem:spread-local-central-hull}
$R$ を可換 Noetherian 環、$\mathfrak p\in\operatorname{Spec}R$ とする。
$B=T_R(V)/(J)$ は連結二次 $R$-代数であり、
$V,J$ はそれぞれ階数四、七の有限 projective 加群、
$J\subset V^{\otimes2}$ は直和因子とする。
$B$ は相対 Koszul であり、全ての $B_n$ は階数 $(n+1)^2$ の
有限 projective 加群として任意の底変更と両立すると仮定する。
$E=B^!$ と置く。

局所環 $R_{\mathfrak p}$ 上で、中央正則元
$\phi_{\mathfrak p}\in(E_{\mathfrak p})_2$ があり、
$P_{\mathfrak p}=E_{\mathfrak p}/(\phi_{\mathfrak p})$ は
階数列 $(1,4,6,4,1)$ の有限自由 Frobenius 代数であるとする。
このとき $\mathfrak p\in D(f)$ となる一つの主開部分へ制限すれば、
$\phi_{\mathfrak p}$ は中央元 $\phi\in(E_f)_2$ に延長し、次が成り立つ。
\begin{enumerate}
\item $\phi$ は全次数で正則であり、この性質は任意の底変更に保たれる。
\item $P=E_f/(\phi)$ は階数列 $(1,4,6,4,1)$ の相対 Koszul
  Frobenius 代数であり、$A=P^!$ の全次数は有限 projective である。
\item $A$ は相対 quadratic Koszul AS4 型の代数であり、
  ある中央正則二次元 $g$ によって $A/(g)\cong B_f$ と書ける。
  この表示、正則性、左右および双加群 Koszul 解消は任意の底変更と両立する。
\end{enumerate}
さらに任意の可換 Noetherian $R_f$-代数 $T$ に対して $B_T$ が左右 Noetherian
ならば、$A_T$ も左右 Noetherian である。
この主張には、$D(f)$ の全ての noncubic fiber において指定した根を
別の ``良い根'' と比較する仮定は要らない。
\end{lemma}

\begin{proof}
全証明は相対モジュライ論文 \cite[noncubic有限型効果性]{QuadricModuli2026} に置く。
\end{proof}

\begin{remark}[一般の Koszul 開性との区別]
上の証明は、指定点の Frobenius 代数 $P_{\mathfrak p}$ が Koszul だから
近傍全体でも Koszul になる、という推論を使っていない。
先に全幾何 fiber の helix を H2 モデルとして認識し、
その同じ even Hom $\mathbb Z$-代数の線形解消から $B_{\bar s}$ の
Koszul 性を得る。この全 fiber の入力と、中央元の一様正則性を経由して
$P_{\bar s}$ の Koszul 性を導く点が重要である。
一般の二次代数では、三次成分の次元を固定したパラメータ空間における
$n$-Koszul locus は開であり、Koszul locus はこれらの可算共通部分である
\cite[Chapter 6, Theorem 2.1 and Corollary 2.2]{PolishchukPositselski2005}。
また同書 Chapter 6, \S6 は、Koszul locus が constructible でない族を与える。
これらを本研究の特別な中央延長に対する否定結果と解するべきではないが、
局所 Koszul 変形原理だけから単一の有限型開近傍での Koszul 性を
推論できるわけではない。本証明は、全幾何 fiber の明示的認識という
別の有限型入力を使って、この問題を回避する。
\end{remark}

\subsection{noncubic 点の有限型近傍での全相対認識}

\begin{theorem}[noncubic 境界を横断する有限型効果性]
\label{thm:noncubic-finite-type-effectivity}
任意の noncubic $\HH$ 対象の normal foundation に対応する点 $x$ と、
定理~\ref{thm:noncubic-hh-trivial-root} の良い根を取る。
有限 foundation の smooth affine chart $(S,x)$ の有限型 affine
\texorpdfstring{\'etale}{etale} 近傍 $(U,u)\to(S,x)$ 上に、
構造対象だけを指定した実際の相対族
\[
 q_U:U\longrightarrow\mathfrak H_2
\]
が存在する。この族の canonical compact 圏は普遍 foundation 代数
$\Lambda_U$ の perfect 圏であり、四項窓は普遍 projective に対応する。

同じ $U$ 上で、選んだ実際の Morita 根 $\mathsf s$ は heart の exact 自己同値となり、
全軌道環 $B$ は二次・相対 Koszul・全次数有限 projective である。
任意の Noetherian 基底変更後に $B$ は左右 Noetherian であり、
同じ根に対する中央正則二次被覆
\[
 A/(g)\simeq B,
 \qquad \rank A_n=\binom{n+3}{3}
\]
を持つ。$A$ は相対 quadratic Koszul AS4 型で、同じ普遍 Noetherian 性を持つ。
全 Hom、canonical tails の比較、相対 Serre 平方同型、被覆の正則性と
左右・双加群 Koszul 解消は任意の導来基底変更と整合する。
根・根同型・被覆はこの族を構成する補助選択であり、$q_U$ の標識には含めない。
\end{theorem}

\begin{proof}
全証明は相対モジュライ論文 \cite[noncubic有限型効果性]{QuadricModuli2026} に置く。
\end{proof}
